\documentclass[11pt,a4paper]{article}
\usepackage[T1]{fontenc}
\usepackage[utf8]{inputenc}
\usepackage[english]{babel}
\usepackage{amsmath,amssymb}
\usepackage[margin=2cm]{geometry}
\setlength{\parindent}{0pt}
\setlength{\emergencystretch}{3em}
\allowdisplaybreaks

\title{About \texttt{worked\_inductions.tex}}
\author{}
\date{}
\begin{document}
\maketitle
This note explains what the test data in \texttt{worked\_inductions.tex} is used for. That file is input for the engine in \texttt{main.c}; it is read by the test \texttt{test\_worked\_inductions.c} in this folder.

\begin{enumerate}
\item \textbf{Purpose.} The file holds 8 equalities in $n$, all true and all proved by induction from $n = 1$. That they are proved is checked elsewhere. Here the engine must also record what makes the proof readable: the two sides computed for the first six values of $n$, the base case, and the inductive step as a computation.

\item \textbf{What it exercises.} Paragraph 40b of \texttt{main.c} and the sentences of chapter 41 of \texttt{output.c}. Every recorded formula is read back from the text that will be printed and checked with numbers, as the formulas of a proof by induction are; in every row of first values the two sides must give the same number; the graph of those values is drawn from them.

\item \textbf{Variety.} Sums whose closed form is a polynomial, with the step worked out as a chain (split off the last term, use the hypothesis, carry out the operations); a sum of powers of $2$; a statement with the sum on the right; and three statements whose step is not worked out, which must say so: a product that gives an unreduced fraction, a sum whose closed form is a fraction, a sum inside a larger expression.

\item \textbf{Expected result.} The line \texttt{PASS: 150 of 150 checks on worked inductions passed}. The expected values were computed by hand; the expected lines of the step are formulas, compared with the recorded ones by value.

\item \textbf{How to run it.} From the top folder of the repository:
\begin{verbatim}
cc -std=c11 -O2 -Wall -Wextra -o /tmp/test_worked_inductions \
    tests/sums_products_inequalities/test_worked_inductions.c \
    && /tmp/test_worked_inductions
\end{verbatim}
\end{enumerate}
\end{document}
